\par\bigskip\noindent\textbf{K101 supporting proof: The complete compact normal stars and all intersections}\par\nobreak
This proof completes the fourteen exceptional/refinement surface stars and
the twelve remaining source-graph curves in each of the two O/T copies.
It uses the genuine point-line coordinates of EI and IA and the common
factor refinement FO. The four strict source surfaces were computed in
SS. The moving selected/reference surface is retained separately.
The characteristic numbers below are the specified fixed-inner
contributions; this note assigns no total parameter return.

\par

\medskip\noindent\textbf{CS.1. The other source graph chart and its genuine line}

\par\nobreak

Use the quotient torus characters
\[
 (T_0,T_1,T_2,T_3)=(\chi^4,\chi\eta,v,V),
 \qquad \sigma=\eta v.
 \tag{CS.1}
\]
At O, \(\eta=-\bar X/\bar Y\) is the actual formal source
coordinate. At T, \(\eta=\bar Y\), and the source equation
solves for \(\bar X\) because its \(\bar X\)-derivative is a unit.
These are precisely the two source graphs already used in EI.
The complementary blowup chart is
\[
 \eta=\sigma p,\qquad
 p=v^{-1},\qquad
 \chi\sigma=T_1T_2.
 \tag{CS.2}
\]
Its inverse on the overlap is \(v=p^{-1}\), \(\eta=\sigma p\).
The source functions are the original functions composed with
these two displayed graph maps. In particular this chart includes
\(p=0\) as a source-graph point.

For every ordinary line-chain label \(j\), take the chart
\[
 X_jY_j=\sigma,\qquad V=\sigma^jX_j,
 \qquad (\chi,p,X_j,Y_j).
 \tag{CS.3}
\]
The residual source quotient and physical circle are literally
\[
 \gamma(\chi,p,X_j,Y_j)
   =(i\chi,p,i^jX_j,i^{-j-1}Y_j),\qquad
 \delta_\theta(X_j,Y_j)
   =(e^{i\theta}X_j,e^{-i\theta}Y_j).
 \tag{CS.4}
\]
Thus CS.2 preserves both actions, including every \(j\bmod4\)
character.

Here \(V\) is the coefficient in the full genuine point-line frame.
To spell this out at O, write
\[
 c_a(\chi,\sigma)=\frac{\bar\tau(aG)}{\sigma},
 \qquad
 r_a=\frac{p-c_a}{p},\qquad
 e_a=\frac{\sigma_a}{r_a}.
 \tag{CS.5}
\]
The two section ideals are \(p\) and \(p-c_a\). The numerator
\(\sigma_a\) is the actual point-line section with exactly their
zero and pole, as in EI\textquotesingle s horizontal Cartier calculation. Their
quotient in CS.5 is therefore a nonzero local line frame through
both section charts. This is a module map on those two Cartier
generators. It remains valid at their coincidence, with the full
source graph retained. In the \(v\)-chart the same function is
\(r_a=1-c_av\); the two formulas are equal on the overlap.
At T the canonical point section gives the same line frame,
since \(aG\) lies in the O chart. Consequently the coefficient
\(V\) is unchanged in CS.2. The \(\mu_4\) action fixes the point
coefficient: \(\bar\tau\) and \(\bar\tau(aG)\) have the same
character, so their ratio in CS.5 is invariant.

The full source reset is the actual odd-\(b\) translation graph.
On the central graph it sends O:\(p\) to T:\(p-b\), and T:\(p\)
to O:\(p-b\), with the full analytic source formula off the
central graph. Its normal line map is the retained Hom unit,
and its line-chain label changes by four. The inverse is the
inverse source graph and reciprocal line unit. The changes of
the source coordinate may contain normal products; their
linearization on a fixed component is the displayed source
map, and differentiation of a nonzero line unit multiplies a
vanishing normal coefficient. Thus the normal map on the
component is precisely the full coefficient map used below.
The period arrows preserve the graph centres and lift to the
same common refinements by their recorded monomial maps.

\par

\medskip\noindent\textbf{CS.2. The finite source-base arrow and the selected value}

\par\nobreak

The native finite source is carried by the literal equations
\[
 t_e=\chi^{-4},\qquad
 x_e=\chi^{-2}\bar X,\qquad
 y_e=\chi^{-3}\bar Y.
 \tag{CS.6}
\]
At the O and T sections the regular source coordinate obeys
\(\eta_e=\chi\eta\). Since \(\epsilon=\chi\sigma\), the
source-graph coordinate \(p=\eta/\sigma=\eta_e/\epsilon\)
is unchanged. The native chain coefficients have the full
arrow
\[
 X_j=\chi^jA_j,\qquad
 Y_j=\chi^{-j-1}B_j,\qquad
 A_jB_j=\epsilon.
 \tag{CS.7}
\]
Both genuine point frames coincide: the numerator and
denominator of \(r_a\) are multiplied by the same \(\chi\)
in CS.6. Thus \(V\) is the same genuine coefficient on the
two sides. The native torus characters are exactly
\[
 (t_e,\eta_e,v,V)=(T_0^{-1},T_1,T_2,T_3).
 \tag{CS.8}
\]
IA supplies their two source charts through the native I2
fibre. The full source equations and the retained point-line
Cartier map extend there, so CS.8 extends over that finite
base value with those actual charts.

At the selected value, use the 39-chart extension of FO.
Its source equation is \(xy=q\), and its line equation is
\[
 UV=\sigma^4q.
 \tag{CS.9}
\]
The ordinary/selected line arrow is
\[
 X_k=\psi^{-1}q^{j_{\mathrm{per}}}s_r,\qquad
 Y_k=\psi q^{-j_{\mathrm{per}}}t_r,\qquad
 k=4j_{\mathrm{per}}+r-1.
 \tag{CS.10}
\]
Here \(j_{\mathrm{per}}\) is the actual period-block label.
It is distinct from the four internal chain-node residues
\(j=0,1,2,3\) in the surface lists. Every component used in
this note is represented in block \(j_{\mathrm{per}}=0\).
This is the block of the actual Q0 extension in FO, rather
than a suppression of an annular integer.

On the central source the chosen Tate representative of
the two-torsion section is the holomorphic unit branch
with limit \(-1\). Its square is \(Q^h\); its valuation at
the selected node is zero, so \(h=0\) and \(P_G=-1\)
on the whole selected disc. For both admitted even values
of \(a\), \(P_G^a=1\) in this representative. The point
generator of CS.5 and its theta generator consequently
have ratio \(\psi=1\) on every component in this note.
This statement includes their section charts by the
Cartier module maps, with no source-period reindex.
For the exceptional p-chart, this restriction can
also be checked in the full coefficients without
evaluating a meromorphic quotient at a zero. Put
\(b=P_G^a\), \(\tau=\sigma p\),
\(\tau_a=\sigma c_a\), and write
\[
 Z(\tau)-1=\tau u(\tau),\qquad
 b-1=\tau_a u(\tau_a).
\]
With \(E\) the full remaining theta-product unit,
the actual quotient is
\[
 \frac{f_b(Z)}{r_a}
 =\frac{Z(\tau)-Z(\tau_a)}{\tau-\tau_a}
   \frac{\tau}{Z(\tau)-1}\frac{E}{b}.
 \tag{CS.10a}
\]
The first factor is the holomorphic divided
difference. At \(\sigma=0\), it cancels the
second factor, and \(E/b=1\). This proves
\(\psi=1\) through \(p=0\) and \(p=c_a\) in
the same nonzero Cartier frame, with its inverse.
Therefore CS.10 has zero extra q exponent and unit one
on these components. The Q0 tangent arrow in the complete
39-chart fan supplies the actual extension at \(q=0\).
It identifies the finite-base normal bundle used below;
there is no unrecorded selected-fibre divisor.

The original finite \(m\)-cap lies on its actual large-\(\sigma\)
domain. These compact components lie in \(\sigma=0\), so
they do not meet that cap. On the surrounding overlap the
ordinary/root square is still
\(\Phi_\beta=A^kB^{j_{\alpha\beta}}\Phi_\alpha\).
In particular the quarter-turn is \(AB^{j_\gamma}\) and
the fourfold word is \(A^4B^{\sum j_\gamma}\), with
\(A^m=B^n\). None of these integers or finite characters
is changed by CS.2, CS.8 or CS.10.

\par

\medskip\noindent\textbf{CS.3. The complete chart stars}

\par\nobreak

All rays below are expressed in the quotient lattice dual
to CS.1. The two added actual boundary rays are
\[
 R_p=(0,1,-1,0),\qquad F_-=(-1,0,0,0).
 \tag{CS.11}
\]
Indeed \(p=T_2^{-1}\) has value one on \(R_p\), while
\(\chi^4,\chi\sigma,V\) have value zero. The finite
coordinate \(t_e=T_0^{-1}\) has value one on \(F_-\),
while \(\eta_e,v,V\) have value zero. These are the
coordinate morphisms CS.2 and CS.8, with their stated
inverses.

The 238-cone ordinary quotient refinement has sixteen
three-faces on which the \(v\)-coordinate is a unit.
Cone each of those faces with \(R_p\). This gives the
sixteen actual p-graph charts. Cone each of the 38
common three-cones with \(F_-\), and add the four
charts
\[
 \langle F_-,R_p,A_j,A_{j+1}\rangle,\qquad
 A_j=(0,1,0,j),\quad 0\le j\le3.
 \tag{CS.12}
\]
The resulting list has \(238+16+38+4=296\) charts.
Every matrix has determinant of absolute value one in
the quotient lattice. The new charts carry CS.4 and
the full source and Hom units of their parent cones.
On overlaps the full coordinate map is exactly
\[
 z_{\rm target}
   =z_{\rm source}^{\,N_{\rm target}^{-1}N_{\rm source}}
       \cdot g_{\rm target,source},
 \tag{CS.13}
\]
where the exponent matrix and every parent source/Hom
unit exponent are stored. The inverse matrix gives
the inverse map on the stated localized overlap.
The coefficients \(g\) are the same actual analytic
units from EI, IA and FO; CS.5 and CS.10 identify
their restrictions on the compact components.

Here is a direct completeness check for each surface
star, rather than a claim about the support of an
unrelated global four-fan. For a fixed face \(F\) of
rank two, choose one incident regular chart. Let
\(Q\) be the two tangent rows of its inverse matrix,
and \(R\) the two tangent columns of its matrix.
Then \(QR=I_2\) and \(Q\) kills both rays of \(F\).
For every incident chart, project its other two rays
by \(Q\). The exact output consists of a cyclic
list \(r_0,\ldots,r_{s-1}\) satisfying
\[
 \det(r_i,r_{i+1})=1,
 \qquad
 \{\text{incident cones}\}
    =\{\langle r_i,r_{i+1}\rangle: i\bmod s\}.
 \tag{CS.14}
\]
Their angular order is checked using integral cross
products and the two half-planes. There are no repeated
cones, and this order goes once around the whole plane.
Thus these actual source charts cover a compact
surface, with exactly the displayed common edges.
The morphisms CS.13 are its full normal overlaps.

For each of the twelve remaining v-unit curves, there
are exactly two incident fixed-curve charts: one old
v-chart and one added p-chart. Their tangent
coordinates are exactly \(v,p\) with \(p=1/v\).
Their full normal rows in CS.13 identify the three
normal lines. This proves that each curve closes to
one \(\mathbb P^1\), with no additional fixed branch.
All fourteen surfaces and all twelve curves are
extracted by their minimal fixed faces, so a shared
chart does not duplicate a component.

\par

\medskip\noindent\textbf{CS.4. Period orbits and the two source copies}

\par\nobreak

The chronological period acts on the labelled source
copies by
\[
 B:(C,j)\longmapsto(\bar C,j+4),\qquad
 C\in\{O,T\},\qquad \bar O=T,\quad\bar T=O.
 \tag{CS.15}
\]
The root action \(\gamma\) fixes O and T; it has
already been taken in the quotient lattice used in
CS.11. For each residue \(j\bmod4\), the orbit of
\((O,j)\) contains \((T,j+4)\) and \((O,j+8)\),
while the orbit of \((T,j)\) contains \((O,j+4)\)
and \((T,j+8)\). These are two distinct orbits.
A power of B which preserves the chain label is zero;
the remaining finite root action does not exchange
the two sections. The source reset \(p\mapsto p-b\)
is the actual map within the image graph and does
not remove the retained chain index.

Hence every local exceptional/refinement type and
every compact p-curve below appears in two source
copies. The four strict source surfaces of SS each
contain both an O and a T section; that geometric
incidence does not identify the two separate
exceptional stars. The actual maps CS.15 determine
their multiplicity.

\par

\medskip\noindent\textbf{CS.5. Normal intersections from the full coefficient maps}

\par\nobreak

On any fixed surface the physical normal weights are
\((-1,+1)\). Let \(m_-\) and \(m_+\) be the corresponding
rows of \(N^{-1}\). The normal character difference is
\(m_D=m_+-m_-\). Subtract the same reference character
\(m_0\) on every chart. Its difference annihilates the
fixed face, so it descends to the tangent quotient:
\[
 \ell_N=(m_D-m_0)R.
 \tag{CS.16}
\]
For a tangent ray \(r\), the Cartier coefficient is
\(\ell_N(r)\). Indeed a coefficient arrow
\(n_{\rm target}=g\,n_{\rm ref}\) gives the vector-frame
identity \(e_{\rm ref}=g\,e_{\rm target}\). The
meromorphic reference vector frame has local
coefficient \(g\), with the positive order of that
function. Applying this to the quotient of the
positive and negative frames gives exactly
\(\ell_N(r)\) for \(D=N_+\otimes N_-^*\).
The two adjacent charts give the same coefficient
by the two normal rows of CS.13; this equality is
checked for every ray in the certificate.
Adding the reference character back is a single
global meromorphic frame; it changes the displayed
divisor by a principal divisor and leaves its square
unchanged.

The integral tangent intersection matrix has entries
\[
 I_{i,i+1}=1,\qquad
 I_{ii}=-\det(r_{i-1},r_{i+1}),\qquad
 I_{ik}=0\quad(k\ne i,i\pm1).
 \tag{CS.17}
\]
The wall relation gives CS.17 directly: the normal
coordinate transverse to the i-th invariant curve
has transition of degree
\(-\det(r_{i-1},r_{i+1})\), and adjacent coordinate
divisors meet once. If \(d\) is the integral
coefficient vector from CS.16, then
\[
 D^2=d^tId,\qquad
 c_1(TF)^2=\mathbf1^tI\mathbf1,\qquad
 \int_Fp_1=\mathbf1^tI\mathbf1-2s.
 \tag{CS.18}
\]
For the tangent formula, the logarithmic coordinate
vector fields give the usual divisor exact sequence
on these explicitly glued toric charts. The top
Euler number is the number \(s\) of their zero-dimensional
orbits, equivalently the zero count of a generic
torus vector field. This proves the last equality
using the same tangent atlas.

As an independent verification, the two actual
auxiliary cocharacters \((2,7,11,23)\) and
\((3,11,17,41)\) give at each chart its tangent weights
\(t_1,t_2\) and normal weights \(n_-,n_+\).
The sums
\[
 \sum\frac{t_1^2+t_2^2}{t_1t_2},
 \qquad
 \sum\frac{(n_+-n_-)^2}{t_1t_2}
 \tag{CS.19}
\]
agree with CS.18 for both choices. Thus the numerical
calculation has a full integral intersection proof
and two separate weight checks. The physical normal
weights themselves remain \((-1,+1)\).

The four exceptional-source surfaces have the
complete tangent rays
\[
 (1,0),(0,1),(-1,0),(0,-1).
 \tag{CS.20}
\]
Their source charts are therefore the explicit four
product charts in coordinates \((q_\infty,v)\),
\((q_\infty,p)\), \((t_e,v)\), \((t_e,p)\), with
\(t_e=q_\infty^{-1}\) and \(p=v^{-1}\).
In negative/positive normal order their nontrivial
coefficient arrows are
\[
 n_{-,{\rm finite}}=q_\infty n_{-,\infty},
 \qquad
 n_{-,p}=v^{j+1}n_{-,v},\qquad
 n_{+,{\rm target}}=n_{+,{\rm source}}.
 \tag{CS.21}
\]
These follow from the full four-by-four matrices,
including FL.8\textquotesingle s refinement factor \(v^j\). In the
stated coefficient convention CS.21 identifies
\[
 N_-=\mathcal O(-1,-j-1),\qquad N_+=\mathcal O
 \quad\text{on }\mathbb P^1\times\mathbb P^1.
 \tag{CS.22}
\]
Both product coordinates and both line clutches are
specified by CS.20-\/-21. Therefore
\(\int p_1=0\) and \(D^2=2(j+1)\).

For a fixed surface with weights \((-1,+1)\), the
weighted normal-link formula in the positive-hyperplane
convention is
\[
 {\cal R}_U(F)=\frac{\int_Fp_1-3D^2}{24}.
 \tag{CS.23}
\]
The actual physical line is \(M=U^*\), by K68\textquotesingle s
nonzero-vector line arrow; its characteristic pair
and its displayed expression both change sign.
The four exceptional-source values are consequently
\[
 {\cal R}_U=-\frac14,-\frac12,-\frac34,-1,
 \qquad \sum {\cal R}_U=-\frac52
 \tag{CS.24}
\]
per O or T copy.

The ten refinement surfaces have the following exact
complete-star values. The IDs are the minimal-face
IDs of the retained fixed-component extraction.

\begin{longtable}[]{@{}
  >{\raggedleft\arraybackslash}p{(\columnwidth - 8\tabcolsep) * \real{0.2000}}
  >{\raggedleft\arraybackslash}p{(\columnwidth - 8\tabcolsep) * \real{0.2000}}
  >{\raggedleft\arraybackslash}p{(\columnwidth - 8\tabcolsep) * \real{0.2000}}
  >{\raggedleft\arraybackslash}p{(\columnwidth - 8\tabcolsep) * \real{0.2000}}
  >{\raggedleft\arraybackslash}p{(\columnwidth - 8\tabcolsep) * \real{0.2000}}@{}}
\toprule\noalign{}
\begin{minipage}[b]{\linewidth}\raggedleft
ID
\end{minipage} & \begin{minipage}[b]{\linewidth}\raggedleft
Number of tangent charts
\end{minipage} & \begin{minipage}[b]{\linewidth}\raggedleft
\(\int p_1\)
\end{minipage} & \begin{minipage}[b]{\linewidth}\raggedleft
\(D^2\)
\end{minipage} & \begin{minipage}[b]{\linewidth}\raggedleft
\({\cal R}_U\)
\end{minipage} \\
\midrule\noalign{}
\endhead
\bottomrule\noalign{}
\endlastfoot
86 & 8 & -12 & 0 & \(-1/2\) \\
94 & 8 & -12 & 1 & \(-5/8\) \\
106 & 8 & -12 & -5 & \(1/8\) \\
118 & 8 & -12 & -3 & \(-1/8\) \\
131 & 6 & -6 & 0 & \(-1/4\) \\
138 & 6 & -6 & -4 & \(1/4\) \\
139 & 6 & -6 & 2 & \(-1/2\) \\
142 & 8 & -12 & -4 & \(0\) \\
146 & 8 & -12 & 0 & \(-1/2\) \\
151 & 8 & -12 & -3 & \(-1/8\) \\
\end{longtable}

Their sum is \(-9/4\) per copy. Every row retains its
complete primitive tangent rays, integral intersection
matrix, normal divisor coefficients, all chart
matrices and all adjacent normal rows in the JSON.
For example ID 86 has tangent rays
\[
 (1,0),(0,1),(-1,1),(0,-1),(1,-1),
 (4,-3),(3,-2),(2,-1)
\]
and normal divisor vector \((0,0,1,0,0,0,0,0)\).
These give \(\int p_1=-12,D^2=0\) by CS.17-\/-18.
The other nine rows use the same displayed method
on their explicitly stored different divisor rows.

\par

\medskip\noindent\textbf{CS.6. The twelve compact p-curves}

\par\nobreak

Let \(M=N_p^{-1}N_v\) be the full character matrix
between the two charts of a curve. In the source
chart, \(v\) is the tangent coordinate. The matched
normal ray has coefficient one in its own normal
column and zero in the other two normal columns.
If its tangent entry is \(e_i\), the actual arrow is
\(n_{i,p}=v^{e_i}g_i n_{i,v}\), with the full point
unit retained by CS.5. Its vector frame therefore
has degree \(-e_i\). This derives the sign from the
specified coefficient map.

The resulting weights and degrees, in the same
normal-ray order, are:

\begin{longtable}[]{@{}rllr@{}}
\toprule\noalign{}
ID & Physical weights & Normal degrees & \({\cal R}_U\) \\
\midrule\noalign{}
\endhead
\bottomrule\noalign{}
\endlastfoot
46 & \((1,-2,1)\) & \((0,-2,1)\) & \(1/8\) \\
47 & \((-3,4,-1)\) & \((-2,0,0)\) & \(1/108\) \\
49 & \((1,-3,2)\) & \((0,-3,2)\) & \(1/24\) \\
51 & \((1,-4,3)\) & \((0,-4,3)\) & \(1/48\) \\
58 & \((1,-2,1)\) & \((0,-4,2)\) & \(1/4\) \\
59 & \((-1,2,-1)\) & \((-3,0,0)\) & \(0\) \\
74 & \((3,-4,1)\) & \((6,-12,3)\) & \(-19/144\) \\
75 & \((-1,2,-1)\) & \((-4,0,0)\) & \(0\) \\
79 & \((-1,3,-2)\) & \((-4,0,0)\) & \(2/9\) \\
83 & \((-1,4,-3)\) & \((-4,0,0)\) & \(5/18\) \\
162 & \((-2,3,-1)\) & \((-4,2,0)\) & \(-1/108\) \\
163 & \((2,-3,1)\) & \((0,-6,3)\) & \(-1/18\) \\
\end{longtable}

For weights \(w_i\) and degrees \(d_i\), the precise
weighted curve expression used here is
\[
 {\cal R}_U
 =-\frac{\sum_i
       \{(4-\sum_h w_h^2)/w_i+2w_i\}d_i}
          {24\prod_iw_i}.
 \tag{CS.25}
\]
It is the expansion of the already proved
normal-link formula, retaining the finite weight
denominators. The sum of the twelve rows is
\(3/4\) per copy, or \(-3/4\) for the actual
physical line. Two independent auxiliary
cocharacters give the same normal degrees from
the two fixed endpoints, with exactly opposite
nonzero tangent weights.

\par

\medskip\noindent\textbf{CS.7. Certificate and scope of the completed contribution}

\par\nobreak

The reproducible program is
\path{code/part3_fourfold/e6_selected_infinity_compact_surface_stars.py.}
Its JSON retains all 296 charts, fourteen complete
surface stars, twelve complete curve stars, every
normal overlap, two auxiliary weight calculations,
the actual parent source/Hom unit exponents and
both input hashes. It passes 1,637 exact checks.
The earlier stored negative divisor vector represented
the dual difference line. That complete source, program
and output are preserved under
\path{records/attempts/e6_selected_infinity_compact_normal_stars/dual_D_divisor_label.}
The present positive vector represents \(D=N_+\otimes N_-^*\)
by the explicit coefficient-to-frame equation in CS.5.
All intersection squares and displayed contributions
are unchanged by that corrected line label.
The two inputs are:

\begin{longtable}[]{@{}
  >{\raggedright\arraybackslash}p{(\columnwidth - 2\tabcolsep) * \real{0.5000}}
  >{\raggedright\arraybackslash}p{(\columnwidth - 2\tabcolsep) * \real{0.5000}}@{}}
\toprule\noalign{}
\begin{minipage}[b]{\linewidth}\raggedright
Input
\end{minipage} & \begin{minipage}[b]{\linewidth}\raggedright
SHA256
\end{minipage} \\
\midrule\noalign{}
\endhead
\bottomrule\noalign{}
\endlastfoot
\path{code/part3/e6_selected_infinity_factor_overlay.json} & \path{af3ef63eed9f324f4aebd9332378961d011944de7f256c98de41a74cf87e8861} \\
\path{code/part3_fourfold/e6_fixed_face_orbits.json} & \path{8917efd4732b3b161eba8b825354a3a46cc9e71193bb7787f44d9ee65d1a4eb4} \\
\end{longtable}

For one O/T copy the newly completed values are
\[
 \sum_{\rm exceptional}{\cal R}_U=-\frac52,\qquad
 \sum_{\rm refinement}{\cal R}_U=-\frac94,\qquad
 \sum_{p\text{-curves}}{\cal R}_U=\frac34.
 \tag{CS.26}
\]
The source-period orbit argument CS.15 gives two
copies. The previously proved ten finite-base
interface curves have sum zero in each copy.
The existing interior-point/curve block and these
new closed stars have disjoint minimal-face IDs;
the four strict source surfaces remain those of
SS and are counted once globally.

These statements identify the compact component
domains, source maps, normal clutches and their
fixed-inner values. The moving selected/reference
surface and the final outer-family/physical-line
orientation map remain separate entries of the
assembly. No total return or free integer is
deduced from CS.26 alone.

